% ============================================================
% 第 12 章 正性（Lead 自译）
% 底本印刷 pp.129–138 = PDF p169–178
% ============================================================
\chapter{正性}

\ChapterAbstract{Fock--Goncharov 坐标把 Hitchin 分支描述为特征簇中“正”的那一部分。本章介绍分裂实李群上正性（positivity）这一基本概念的来龙去脉，以及由 Guichard--Wienhard 在 [179, 180] 中引入的、适用于实半单李群的广义正性概念——\emph{Θ-正性}（Θ-positivity），它刻画了高阶 Teichmüller 空间。该概念源于在群 $G$ 对应的李代数 $\mathfrak{g}$ 的全部根中，选取某一子集的单正根 $\Theta$。凑巧的是，希腊字母 $\Theta$ 也正是希腊语单词 $\Theta\varepsilon\tau\iota\kappa\acute{o}\tau\eta\tau\alpha$（意为“正性”）的首字母。本章还给出了群 $\SL_n(\R)$ 与 $\Sp_{2n}(\R)$ 的 Θ-正结构的详细分析。}

% [原文 p.129]
\section{全正性}

全正性（total positivity）这一概念出现在广得出人意料的众多数学领域中。本节我们简要回顾 10.1.2 节给出的\emph{全正矩阵}（totally positive matrix）的定义，并展开讨论其性质。一个元素为实数的可逆 $n\times n$ 矩阵称为\emph{全正的}，如果它的所有方子式（square submatrices）的行列式都为正。下面我们展示该概念的一些重要特征，关于一般性的论述可参见 [20, 163, 227]。

\begin{definition}[12.1]
称矩阵 $m\in\GL_n(\R)$ 是\emph{全非负的}（totally non-negative）（相应地，\emph{全正的}），如果对任意满足 $|I|=|J|$ 的 $I,J\subset[\![1,n]\!]$，都有 $\Delta_{I,J}(m)\in\R_{\ge 0}$（相应地，$\Delta_{I,J}(m)\in\R_{>0}$）；这里 $\Delta_{I,J}(m)$ 表示 $m$ 中由指标集 $I$ 标定的行与由指标集 $J$ 标定的列所构成的子式。
\end{definition}

\begin{example}[12.2]
矩阵
\[
M=\begin{pmatrix} 22 & 9 & 6\\ 8 & 4 & 3\\ 1 & 1 & 1\end{pmatrix}
\]
是全正的，因为它的所有元素、所有 $2\times 2$ 行列式以及 $\det(M)$ 都严格为正。
\end{example}

% [原文 p.130]
\begin{proposition}[12.3]
全非负（相应地，全正）矩阵构成 $\GL_n(\R)$ 的一个乘法子半群，记为 $\GL_{\ge 0}$（相应地，$\GL_{>0}$）。
\end{proposition}

\begin{proof}
这是“两个方阵之积的行列式等于它们行列式之积”这一性质的直接推论。
\end{proof}

% [Tutor Note: 注意这里 $m\in\GL_n(\R)$ 只要求可逆；全非负矩阵的行列式非负，故其行列式只能为正（不能为 0），全正的矩阵当然非奇异。Prop. 12.3 中“半群”而非“群”是因为全非负矩阵的逆不一定全非负。]

为了更好地理解 $\GL_n(\R)$ 中全正矩阵的结构，我们先引入它的几个子群。设 $B$ 为 $\GL_n(\R)$ 中全体上三角矩阵构成的子群，称为\emph{Borel 子群}（Borel subgroup），也称作\emph{极大抛物子群}；再设 $B^{\mathrm{opp}}$ 为 $\GL_n(\R)$ 中全体下三角矩阵构成的子群，称为\emph{相反 Borel 子群}（opposite Borel subgroup）。又设 $U$ 为 $B$ 中全体对角元为 1 的矩阵构成的子群，称为\emph{上单幂子群}（upper unipotent subgroup），或简称为\emph{单幂子群}；并设 $U^{\mathrm{opp}}$ 为 $B^{\mathrm{opp}}$ 中全体对角元为 1 的矩阵构成的子群，称为\emph{下单幂子群}（lower unipotent subgroup）。

全正性的概念可以推广到 $U$ 与 $U^{\mathrm{opp}}$ 中的矩阵，如下：

\begin{definition}[12.4]
称矩阵 $m\in U\cup U^{\mathrm{opp}}$ 是\emph{全非负的}（相应地，\emph{全正的}），如果对任意满足 $|I|=|J|$ 的 $I,J\subset[\![1,n]\!]$，都有 $\Delta_{I,J}(m)\in\R_{\ge 0}$（相应地，$\Delta_{I,J}(m)\in\R_{>0}$），但那些因 $m\in U$（相应地，$m\in U^{\mathrm{opp}}$）而必须为零的子式除外。$U$（相应地，$U^{\mathrm{opp}}$）的全体全非负（相应地，全正）矩阵之集记为 $U_{\ge 0}$（相应地，$U_{>0}$）以及 $U^{\mathrm{opp}}_{\ge 0}$（相应地，$U^{\mathrm{opp}}_{>0}$）。
\end{definition}

\begin{example}[12.5]
矩阵 $U_1=\begin{pmatrix} 1 & 3 & 6\\ 0 & 1 & 3\\ 0 & 0 & 1\end{pmatrix}$ 是单幂的，并且容易验证它是 $U$ 中的一个全正元素。
\end{example}

下述性质是研究全正矩阵的一个关键事实：

\begin{proposition}[12.6（Cryer 分裂引理，Cryer's splitting lemma）\ {{{{{[97, 98]}}}}}]
可逆矩阵 $m\in\GL_n(\R)$ 是全非负的，当且仅当它有一个 Gauss 分解 $M=U_1DU_2$，其中 $U_1\in U_{\ge 0}$，$D$ 是对角元非负的对角矩阵，且 $U_2\in U^{\mathrm{opp}}_{\ge 0}$。
\end{proposition}

\begin{example}[12.7]
回到例 12.2，矩阵 $M$ 的 Gauss 分解为
\[
M=\begin{pmatrix} 22 & 9 & 6\\ 8 & 4 & 3\\ 1 & 1 & 1\end{pmatrix}
=\begin{pmatrix} 1 & 3 & 6\\ 0 & 1 & 3\\ 0 & 0 & 1\end{pmatrix}
\begin{pmatrix} 1 & 0 & 0\\ 5 & 1 & 0\\ 1 & 1 & 1\end{pmatrix}.
\]
对角部分 $D$ 是平凡（单位矩阵），$U_1$ 由例 12.5 知是全正的，而 $U_2$ 也容易验证是全正的。
\end{example}

% [Tutor Note: 原文说“the Gauss decomposition ... is given by”，但只写出两个因子；按 Prop. 12.6 的记号应为 $M=U_1DU_2$，此处 $D=I_3$ 已被略去。验算：$\begin{pmatrix}1&3&6\\0&1&3\\0&0&1\end{pmatrix}\begin{pmatrix}1&0&0\\5&1&0\\1&1&1\end{pmatrix}$ 的第 1 行 $=(1+15+6,\ 3+6,\ 6)=(22,9,6)$，与 $M$ 相符。]

% [原文 p.131]
随后，对半群 $\GL_{\ge 0}$ 的研究就归结为对其子半群 $U_{\ge 0}\subset\GL_{\ge 0}$（上三角单幂全非负矩阵）的研究。关于全非负矩阵的下述基本约化定理，使我们可以把 $U_{\ge 0}$ 的无穷小生成元等同于相应李代数的 Chevalley 生成元：

\begin{theorem}[12.8（Loewner--Whitney {{{{{[248, 360]}}}}}）]
任意可逆的全非负矩阵都是初等矩阵 $x_i(t)$（见 (12.1)）的乘积，其中 $t>0$。
\end{theorem}

该定理蕴含：每个非负的单幂上三角矩阵 $u\in U_{\ge 0}$ 都可写成全非负矩阵 $x_i(t)$ 的乘积，其中 $t\in\R_{\ge 0}$，而 $x_i(t)$ 是 $\GL_n(\R)$ 中对角元为 1、第 $(i,i+1)$ 元为 $t$、其余元素为 0 的矩阵：
\begin{equation}
x_i(t)=
\begin{pmatrix}
1 & \cdots & 0 & 0 & \cdots & 0\\
 & \ddots & \vdots & \vdots & \\ddots & \\
0 & \cdots & 1 & t & \cdots & 0\\
0 & \cdots & 0 & 1 & \cdots & 0\\
 & \\ddots & \vdots & \vdots & \ddots & \\
0 & \cdots & 0 & 0 & \cdots & 1
\end{pmatrix}.
\tag{12.1}
\end{equation}

\begin{example}[12.9]
回到例 12.5，我们容易求出全正矩阵 $U_1$ 的分解：
\[
U_1=\begin{pmatrix} 1 & 3 & 6\\ 0 & 1 & 3\\ 0 & 0 & 1\end{pmatrix}
=\begin{pmatrix} 1 & 2 & 0\\ 0 & 1 & 0\\ 0 & 0 & 1\end{pmatrix}
\begin{pmatrix} 1 & 0 & 0\\ 0 & 1 & 3\\ 0 & 0 & 1\end{pmatrix}
\begin{pmatrix} 1 & 1 & 0\\ 0 & 1 & 0\\ 0 & 0 & 1\end{pmatrix}
=x_1(2)x_2(3)x_1(1).
\]
\end{example}

这一性质提供了把全正性概念推广到 $\GL_n(\R)$ 之外的想法。具体地说，Lusztig 在 [251] 中把分裂实约化李群 $G$ 中全非负元素之集 $G_{\ge 0}$ 定义为由相应的 Chevalley 生成元生成的半群。其主要想法是考虑把李代数的单根空间取指数所得的单参数子群，而坐标变换则由正有理映射给出。Borel 子群与单幂子群的概念可以相应地推广到任意分裂实李群，见 A.3 节。

分裂实群 $G$ 中全正元素构成的子半群，还与广义旗簇 $G/B$ 中三元组（triples）的正性概念密切相关，这里 $B$ 是 $G$ 的 Borel 子群。$G/B$ 的元素可以自然地等同于 $G$ 中那些在 $G$ 下共轭于标准 Borel 子群 $B$ 的子群。旗对 $(B,B^{\mathrm{opp}})$（其中 $B^{\mathrm{opp}}$ 是 $G$ 的相反 Borel 子群）称为\emph{标准横截对}（standard transverse pair）。两个旗 $E$ 与 $F$ 称为\emph{横截的}（transverse），如果存在 $g\in G$ 使得 $(gE,gF)=(B,B^{\mathrm{opp}})$。

\begin{definition}[12.10]
设 $G$ 是分裂实李群。称旗的三元组 $(E,T,F)$ 是\emph{正的}，当且仅当存在 $g\in G$ 与 $u\in U_{>0}$ 使得
\[
(gE,gT,gF)=(B,uB^{\mathrm{opp}},B^{\mathrm{opp}}).
\]
\end{definition}

\begin{example}[12.11]
对 $G=\GL_3(\R)$，一个旗可以表示为 $\R P^2$ 中的一条直线连同其上一个标定点。当三比值（triple ratio）为正时，称相应的旗三元组为正。
\end{example}

% [Tutor Note: 原文此句在印刷本中即不完整：“A triple of flags is positive if the triple ratio is.” 句末缺宾语（推测应为“is positive”）。此处按上下文补全为“当三比值……时”，并以批注标记；底本疑有脱字，参校时请以 Guichard–Wienhard 原文为准。]

% [原文 p.132]
\section{Θ-正表示}

Θ-正表示（Θ-positive representation）这一概念是 Lusztig 在 [251] 中提出的全正性条件的推广，适用于未必分裂的实半单李群。与此同时，它也推广了 [70] 中所阐述的 Hermitian 李群的极大性（maximality）的那些重要性质；而它在高阶 Teichmüller 理论中的意义，从 [179] 或 [177] 的结果中已经显而易见。本节我们介绍这一概念，并在以下各小节给出详细例子；关于 Θ-正表示的进一步参考，读者可查阅原始文献 [179]、[180]、[177]。在接下来的几节中会给出显式计算，这对熟悉以下对象大有裨益。

关于李论的背景（包括 Cartan 子代数、根空间、Levi 子代数与抛物子代数），可参见附录 A。实半单李群 $G$ 的 Θ-正结构的定义，是借助由单正根的子集 $\Theta\subset\Delta$ 所定义的抛物子群 $P_\Theta<G$ 的李代数的性质给出的。按这些记号，我们首先考虑中心有限的实半单李群 $G$，并记 $\mathfrak{g}$ 为 $G$ 的李代数。设 $K$ 为 $G$ 的极大紧子群，$\mathfrak{k}$ 为其李代数，$\mathfrak{k}^\perp$ 为 $\mathfrak{k}$ 关于 Killing 型的正交补。进而设 $\mathfrak{a}$ 为 $\mathfrak{k}^\perp$ 的一个极大交换 Cartan 子空间，并记 $\Sigma=\Sigma(\mathfrak{g},\mathfrak{a})$ 为限制根系（system of restricted roots）。最后，取 $\Delta\subset\Sigma$ 为一组单根，并记 $\Sigma^+$ 为关于 $\Delta$ 的正根之集。

现在，设 $\Theta\subset\Delta$ 为单根的任意子集（$\Theta$ 可以为空）。令 $\Sigma_\Theta^+=\Sigma^+\setminus\operatorname{Span}(\Delta\setminus\Theta)$，并令
\[
\mathfrak{u}_\Theta=\bigoplus_{\alpha\in\Sigma_\Theta^+}\mathfrak{g}_\alpha,
\qquad
\mathfrak{u}_\Theta^{\mathrm{opp}}=\bigoplus_{\alpha\in\Sigma_\Theta^+}\mathfrak{g}_{-\alpha}
\]
以及
\[
\mathfrak{l}_\Theta=\mathfrak{g}_0\oplus\bigoplus_{\alpha\in\Sigma^+\cap\operatorname{Span}(\Delta\setminus\Theta)}\bigl(\mathfrak{g}_\alpha\oplus\mathfrak{g}_{-\alpha}\bigr).
\]
作为向量空间，$\mathfrak{g}$ 分解为 Lie 子代数的直和 $\mathfrak{g}=\mathfrak{u}_\Theta^{\mathrm{opp}}\oplus\mathfrak{l}_\Theta\oplus\mathfrak{u}_\Theta$。此外，$\mathfrak{l}_\Theta$ 通过 Lie 括号作用在 $\mathfrak{u}_\Theta$ 上；我们接下来要理解这一作用及其性质。

% [Tutor Note: 这里的 $\Sigma_\Theta^+=\Sigma^+\setminus\operatorname{Span}(\Delta\setminus\Theta)$ 是集合差，不是线性空间的商或补。于是 $\Sigma_\Theta^+$ 由那些“不能只用 $\Delta\setminus\Theta$ 中单根的整系数线性组合表示”的正根组成；若 $\Theta=\Delta$，则 $\Sigma_\Theta^+=\Sigma^+$；若 $\Theta=\varnothing$，则 $\Sigma_\Theta^+=\varnothing$，此时 $\mathfrak{u}_\Theta=0$、$\mathfrak{l}_\Theta=\mathfrak{g}$。]

群 $G$ 通过伴随表示作用在 $\mathfrak{g}$ 上；设 $P_\Theta$ 为 $\mathfrak{u}_\Theta$ 在 $G$ 中的正规化子。$P_\Theta$ 的李代数为 $\mathfrak{p}_\Theta=\mathfrak{l}_\Theta\oplus\mathfrak{u}_\Theta$。类似地，记 $P_\Theta^{\mathrm{opp}}$ 为 $\mathfrak{u}_\Theta^{\mathrm{opp}}$ 的正规化子，其李代数为 $\mathfrak{p}_\Theta^{\mathrm{opp}}=\mathfrak{l}_\Theta\oplus\mathfrak{u}_\Theta^{\mathrm{opp}}$。这些就是与 $\Theta$ 相联系的标准抛物子群，相应的 Levi 子群为 $L_\Theta=P_\Theta\cap P_\Theta^{\mathrm{opp}}$。子群 $U_\Theta=\exp(\mathfrak{u}_\Theta)$（相应地，$U_\Theta^{\mathrm{opp}}=\exp(\mathfrak{u}_\Theta^{\mathrm{opp}})$）是 $P_\Theta$（相应地，$P_\Theta^{\mathrm{opp}}$）的单幂根（unipotent radical）。$L_\Theta$ 的李代数是 $\mathfrak{l}_\Theta$，但 $L_\Theta$ 未必连通，故记 $L_\Theta^0$ 为 $L_\Theta$ 中含单位元的连通分支。

为理解 $L_\Theta$ 在 $\mathfrak{u}_\Theta$ 上的作用，由于 $L_\Theta$ 是约化群，我们首先需要理解其中心与 Cartan 子代数。首先，$\mathfrak{l}_\Theta$ 分解为中心 $Z(\mathfrak{l}_\Theta)$ 与其半单部分 $\mathfrak{l}'_\Theta=[\mathfrak{l}_\Theta,\mathfrak{l}_\Theta]$ 的直和：
\[
\mathfrak{l}_\Theta=Z(\mathfrak{l}_\Theta)\oplus\mathfrak{l}'_\Theta.
\]

% [原文 p.133]
子群 $Z(L_\Theta)$ 的李代数 $\mathfrak{z}_\Theta$ 再一次分成两部分：$\mathfrak{z}_\Theta=Z(\mathfrak{l}_\Theta)\cap\mathfrak{k}^\perp=Z(\mathfrak{l}_\Theta)\cap\mathfrak{a}$ 以及 $Z(\mathfrak{l}_\Theta)\cap\mathfrak{k}$。

$\mathfrak{l}'_\Theta$ 的 Cartan 子代数为 $\mathfrak{a}_\Theta=\mathfrak{a}\cap\mathfrak{l}'_\Theta$，其限制根系是 $\Sigma\cap\operatorname{Span}(\Delta\setminus\Theta)$ 到 $\mathfrak{a}_\Theta$ 的限制，一组可能的单根是 $\Delta\setminus\Theta$ 到 $\mathfrak{a}_\Theta$ 的限制，而相应的 Dynkin 图是 $G$ 的 Dynkin 图中以 $\Delta\setminus\Theta$ 为顶点的最大子图。

注意，由于 $\mathfrak{a}\subset\mathfrak{l}_\Theta\cap\mathfrak{k}^\perp$，我们有
\[
\mathfrak{a}=\mathfrak{a}_\Theta\oplus\mathfrak{z}_\Theta.
\]

特别地，$\mathfrak{z}_\Theta$ 通过 $\mathfrak{a}$ 在 $\mathfrak{g}$ 上作用的限制而作用在 $\mathfrak{u}_\Theta$ 上，因此 $\mathfrak{u}_\Theta$ 按此作用的权空间分解为
\begin{equation}
\mathfrak{u}_\Theta=\bigoplus_{\beta\in\mathfrak{z}_\Theta^*}\mathfrak{u}_\beta,
\tag{12.2}
\end{equation}
其中
\begin{equation}
\mathfrak{u}_\beta=\bigl\{x\in\mathfrak{u}_\Theta\mid\ \text{对一切 } z\in\mathfrak{z}_\Theta,\ \operatorname{ad}(z).x=\beta(z)x\bigr\}
\tag{12.3}
\end{equation}
是 $L_\Theta^0$-不变的。

包含映射 $\mathfrak{z}_\Theta\subset\mathfrak{a}$ 诱导出一个投影 $i^*:\mathfrak{a}^*\to\mathfrak{z}_\Theta^*$，它由线性形式在 $\mathfrak{a}$ 上限制到 $\mathfrak{z}_\Theta$ 给出。于是我们可以把 $\mathfrak{u}_\beta$ 写成 $\mathfrak{a}$ 的作用的权空间之和：
\[
\mathfrak{u}_\beta=\bigoplus_{\alpha\in\Sigma_\Theta^+,\ \alpha|_{\mathfrak{z}_\Theta}=\beta}\mathfrak{g}_\alpha.
\]

事实上，$\mathfrak{u}_\beta$ 是不可约的：

\begin{theorem}[12.12（{{{{{[240, Thm. 0.1]}}}}}）]
对一切 $\beta\in\mathfrak{z}_\Theta^*$，$\mathfrak{u}_\beta$ 都是 $L_\Theta^0$ 的不可约表示，并且
\[
[\mathfrak{u}_\beta,\mathfrak{u}_{\beta'}]=\mathfrak{u}_{\beta+\beta'}.
\]
\end{theorem}

% [Tutor Note: 定理中的 $[\mathfrak{u}_\beta,\mathfrak{u}_{\beta'}]=\mathfrak{u}_{\beta+\beta'}$ 按惯例应理解为：若 $\beta+\beta'$ 是 $\mathfrak{z}_\Theta$ 在 $\mathfrak{u}_\Theta$ 上的一个权，则左端等于相应的 $\mathfrak{u}_{\beta+\beta'}$，否则左端为 $0$（因为 $\mathfrak{u}_\Theta$ 对李括号封闭，括号结果落在 $\mathfrak{u}_\Theta$ 内）。]

设 $\beta$ 是某个根 $\alpha\in\Theta$ 的限制。则 $\alpha$ 是 $\operatorname{Span}(\Theta)$ 中唯一限制为 $\beta$ 的元素，这就允许我们作如下略微滥用的记号：
\begin{equation}
\mathfrak{u}_\alpha:=\mathfrak{u}_\beta=\bigoplus_{\gamma\in\operatorname{Span}_{\N}(\Delta\setminus\Theta)}\mathfrak{g}_{\alpha+\gamma}.
\tag{12.4}
\end{equation}

\begin{definition}[12.13]
设 $V$ 为有限维实向量空间。$V$ 中的\emph{锥}（cone）是指一个对正数乘法封闭的子集 $C$，即对一切 $x\in C$ 与一切 $\lambda>0$ 都有 $\lambda x\in C$。若锥 $C$ 不含仿射直线，则称它是\emph{尖锐的}（sharp）。
\end{definition}

\begin{remark}[12.14]
尖锐性条件保证锥在某种意义下不会“过宽”；例如半平面 $\HH^2\subset\R^2$ 是一个锥，但不尖锐，因为它包含许多水平方向的仿射直线。
\end{remark}

% [Tutor Note: 原文用 $\HH^2$ 表示上半平面（第 4 章约定），此处“半平面”指 $\{y>0\}$ 这一开半平面；它是锥（对正数乘法封闭）但不尖锐（含水平直线）。注意不要把尖锐性与“闭”或“凸”混淆：Def. 12.13 只要求对正数乘法封闭。]

% [原文 p.134]
我们现在终于拥有了给出 Θ-正性定义所需的全部工具与记号；该定义统一并推广了我们在 10.1 节与 11.1 节中分别对分裂实李群与管型 Hermitian 李群所详尽描述的正性概念。

\begin{definition}[12.15（Guichard--Wienhard，{{{{{[179, 180]}}}}}）]
设 $G$ 为中心有限的实半单李群，$\Theta\subset\Delta$ 为单根的子集。若对一切 $\beta\in\Theta$，$L_\Theta^0$ 在 $\mathfrak{u}_\beta$ 上的作用都保持某个尖锐凸锥（sharp convex cone），则称 $G$ 具有一个 \emph{Θ-正结构}（Θ-positive structure）。
\end{definition}

[179, 180] 中的一个核心结果给出了具有 Θ-正结构的半单李群 $G$ 的分类：

\begin{theorem}[12.16（Guichard--Wienhard，{{{{{[180]}}}}} 中 Thm. 3.4）]
半单李群 $G$ 具有 Θ-正结构，当且仅当对 $(G,\Theta)$ 属于下列四种情形之一：
\begin{enumerate}
\item $G$ 是分裂实群且 $\Theta=\Delta$；
\item $G$ 是管型的 Hermitian 对称李群且 $\Theta=\{\alpha_r\}$；
\item $G$ 局部同构于群 $\SO(p,q)$（$p\ne q$），且 $\Theta=\{\alpha_1,\dots,\alpha_{p-1}\}$；
\item $G$ 是群 $F_4$、$E_6$、$E_7$、$E_8$ 的实形式，其限制根系为 $F_4$ 型，且 $\Theta=\{\alpha_1,\alpha_2\}$。
\end{enumerate}
\end{theorem}

% [Tutor Note: 情形 3 中 $p\ne q$ 是必要的：当 $p=q$ 时 $\SO(p,p)$ 是分裂实群，属于情形 1，且 $\Theta=\Delta$ 给出全正性；$\Theta=\{\alpha_1,\dots,\alpha_{p-1}\}$ 这种取法对应把 $\SO(p,q)$ 实现为某个管型 Hermitian 李群的实形式。另注意情形 2 中 $\alpha_r$ 指 Dynkin 图的“末端”根（与紧 Cartan 方向对应的那个单根），本书 12.4 节的 $\Sp_{2n}(\R)$ 例子正是 $r=n$ 的情形。]

人们还可以把广义旗簇 $G/P_\Theta$ 中的正三元组（positive triples）考虑如下：

\begin{definition}[12.17]
固定 $G/P_\Theta$ 中的标准旗 $E_\Theta$ 与 $F_\Theta$，使得 $\operatorname{Stab}_G(F_\Theta)=P_\Theta$ 且 $\operatorname{Stab}_G(E_\Theta)=P_\Theta^{\mathrm{opp}}$。对任意与 $F_\Theta$ 横截的 $S_\Theta\in G/P_\Theta$，存在 $u_{S_\Theta}\subset U_\Theta$ 使得 $S_\Theta=u_{S_\Theta}E_\Theta$。若 $u_{S_\Theta}\in U_\Theta^{>0}$，则称广义旗簇 $G/P_\Theta$ 中的三元组 $(E_\Theta,S_\Theta,F_\Theta)$ 是 \emph{Θ-正的}（Θ-positive）；这里 $U_\Theta^{>0}$ 是 $U_\Theta$ 的 Θ-正子半群，其定义方式与 Lusztig 在分裂实群 $G$ 情形下对正元素之集 $U_{>0}$ 的定义方式类似（确切定义见 [180] 第 4 节）。
\end{definition}

Θ-正基本群表示的定义现在可以给出如下：

\begin{definition}[12.18]
设 $S_g$ 是亏格 $g\ge 2$ 的闭连通可定向拓扑曲面，$G$ 是容许 Θ-正结构的半单李群。若存在一个 $\rho$-等变的正映射 $\xi:\partial\pi_1(S_g)=\R P^1\to G/P_\Theta$，把 $\R P^1$ 中的正三元组送到 $G/P_\Theta$ 中的 Θ-正三元组，则称基本群 $S_g$ 到 $G$ 的一个表示是 \emph{Θ-正的}。
\end{definition}

% [Tutor Note: Def. 12.18 中的 $\partial\pi_1(S_g)=\R P^1$ 指曲面 $S_g$ 的（Gromov）边界，等同为圆 $\R P^1$；$\rho$-等变指 $\xi$ 与 $G$ 在 $G/P_\Theta$ 上的作用以及 $\rho$ 在边界上的作用相容。一个 $\Theta$-正表示必然离散且忠实（见 Thm. 12.19），这正是“正性 $\Rightarrow$ 离散忠实”这一 Hitchin 型现象的一般形式。]

% [原文 p.135]
在 [177] 中，Guichard、Labourie 与 Wienhard 证明了：Θ-正表示是 Θ-Anosov 的，从而是离散且忠实的；并且，对上文 Thm. 12.16 所列的四族半单李群 $G$，特征簇中都存在高阶 Teichmüller 空间：

\begin{theorem}[12.19（Guichard--Labourie--Wienhard，{{{{{[177]}}}}} 中 Thm. A）]
设 $G$ 是容许 Θ-正结构的半单李群。则存在表示簇 $\mathcal{R}(S_g,G)$ 的一个连通分支，它只由离散且忠实的表示组成。
\end{theorem}

\section{$\SL_n(\R)$ 中的 Θ-正性}

下面我们以群 $\SL_n(\R)$ 为例详细展示 Θ-正结构。其李代数是迹为零的矩阵之集。

沿用前面引入的一般记号，现取 $G=\SL_n(\R)$，$K=\SO_n(\R)$ 为 $G$ 的极大紧子群，其李代数 $\mathfrak{k}$ 为斜对称矩阵之集。$\mathfrak{k}$ 关于 Killing 型的正交补 $\mathfrak{k}^\perp$ 是迹为零的对称矩阵之集，而 $\mathfrak{k}^\perp$ 的一个极大交换 Cartan 子空间 $\mathfrak{a}$ 是迹为零的对角矩阵之集。此时限制根系 $\Sigma=\Sigma(\mathfrak{g},\mathfrak{a})$ 为
\[
\Sigma=\bigl\{\alpha_{ij}:=\epsilon_i-\epsilon_j\ \big|\ i,j\in\{1,\dots,n\},\ i\ne j\bigr\},
\]
其中 $\epsilon_i:(a_{pq})_{p,q\in\{1,\dots,n\}}\mapsto a_{ii}$。单根的标准取法为 $\Delta=\{\alpha_{i,i+1}\}_{i\in\{1,\dots,n-1\}}$，而关于 $\Delta$ 的正根之集为 $\Sigma^+=\{\alpha_{ij}\}_{i<j}$。权空间 $\mathfrak{g}_{\alpha_{ij}}$ 是由 $E_{ij}$ 张成的 $\mathfrak{sl}_n(\R)$ 的一维子空间。

\begin{example}[12.20]
我们首先描述一种不给出 Θ-正结构的 $\Theta\subset\Delta$ 的取法，这将有助于我们更好地理解与 $\Theta$ 相联系的各个李代数。取 $\Theta=\{\alpha_{i,i+1},\alpha_{j,j+1}\}$，其中 $1\le i<j<n$。则李代数 $\mathfrak{sl}_n(\R)$ 具有如下分块分解
\[
\left(
\begin{array}{c|c|c}
\ast & \ast & \ast\\ \hline
\ast & \ast & \ast\\ \hline
\ast & \ast & \ast
\end{array}
\right)
\quad\text{其中 } \Tr=0.
\]
\end{example}

% [Tutor Note: 原文此处是一张示意矩阵图（见底本 p.136）：三条粗线把 $n\times n$ 矩阵分成 $3\times 3$ 块，块的大小依次为 $i$、$j-i$、$n-j$，左侧标注行指标 $i,i+1,j,j+1$，右侧写“with $\Tr=0$”。这里用 $3\times3$ 的 $\ast$-块矩阵复现其结构，$\ast$ 表示任意元素。]

% [原文 p.136]
下文中，当我们谈论按块定义的矩阵时，将引用这一分块分解。Levi 子代数 $\mathfrak{l}_\Theta$ 是迹为零的分块对角矩阵之集，代数 $\mathfrak{u}_\Theta$ 是分块严格上三角矩阵之集，而 $\mathfrak{p}_\Theta=\mathfrak{l}_\Theta\oplus\mathfrak{u}_\Theta$ 是迹为零的分块上三角矩阵之集：
\[
\mathfrak{l}_\Theta=\left\{\left(
\begin{array}{c|c|c}
\ast & & \\ \hline
 & \ast & \\ \hline
 & & \ast
\end{array}
\right),\ \Tr=0\right\},\qquad
\mathfrak{u}_\Theta=\left\{\left(
\begin{array}{c|c|c}
 & \ast & \ast\\ \hline
 & & \ast\\ \hline
 & &
\end{array}
\right),\ \Tr=0\right\},
\]
\[
\mathfrak{p}_\Theta=\left\{\left(
\begin{array}{c|c|c}
\ast & \ast & \ast\\ \hline
 & \ast & \ast\\ \hline
 & & \ast
\end{array}
\right),\ \Tr=0\right\}.
\]

抛物子群 $P_\Theta$ 是 $\R^n$ 中以 $(0,i,j,n)$ 为维数序列的标准（部分）旗的稳定化子。

对于我们选取的单根 $\Theta=\{\alpha_{i,i+1},\alpha_{j,j+1}\}$（$1\le i<j<n$），$\mathfrak{l}_\Theta$ 的中心 $Z(\mathfrak{l}_\Theta)$ 是形如
\[
\left(
\begin{array}{c|c|c}
\lambda I_i & & \\ \hline
 & \mu I_{j-i} & \\ \hline
 & & \nu I_{n-j}
\end{array}
\right)
\]
的矩阵之集，且满足 $i\lambda+(j-i)\mu+(n-j)\nu=0$。此时 $\mathfrak{l}_\Theta$ 的半单部分是那些每个块都迹为零的分块对角矩阵之集。在这种情形下，$Z(\mathfrak{l}_\Theta)$ 完全包含在 $\mathfrak{a}$ 中，故 $\mathfrak{z}_\Theta=Z(\mathfrak{l}_\Theta)$。权空间 $\mathfrak{u}_\beta$ 是 $\mathfrak{u}_\Theta$ 中只有三个块之一非零的分块上三角矩阵；而 $\beta\in\Theta$ 对应的 $\mathfrak{u}_\beta$ 就是紧贴对角线之上的那些块。第三个 $\mathfrak{u}_\beta$（其唯一的非零块位于矩阵右上角）对应于根 $\alpha_{i,j+1}=\alpha_{i,i+1}+\cdots+\alpha_{j,j+1}$ 到 $\mathfrak{z}_\Theta$ 的限制，而所有根 $\alpha_{i+1,i+2},\dots,\alpha_{j-1,j}$ 在 $\mathfrak{z}_\Theta$ 上为零。群 $L_\Theta^0$ 是形如

% [原文 p.137]
\[
\left(
\begin{array}{c|c|c}
\lambda A & & \\ \hline
 & \mu B & \\ \hline
 & & \nu C
\end{array}
\right)
\]
的矩阵之集，其中 $A,B,C$ 是行列式为正的可逆矩阵。在这种情形下，$L_\Theta^0$ 在 $\mathfrak{u}_\beta$ 上的作用不保持任何尖锐凸锥。

% [Tutor Note: 这正是例 12.20 的结论：$\Theta=\{\alpha_{i,i+1},\alpha_{j,j+1}\}$（$i$、$j$ 不是相邻指标）不给出 Θ-正结构。直观上，右上角的 $\mathfrak{u}_\beta$ 是 $L_\Theta^0$ 的 1 维表示，但 $\lambda,\mu,\nu$ 受 $i\lambda+(j-i)\mu+(n-j)\nu=0$ 约束而符号不定，故无法取到一个与所有群元素相容的尖锥。]

我们下面说明：全正性如何在单根的某种特定取法下成为更一般的 Θ-正性概念的特例。对群 $G=\SL_n(\R)$，取 $\Theta=\Delta$（$\Delta$ 取 $\SL_n(\R)$ 的通常取法）。此时 Levi 子群 $L_\Theta$ 是行列式为 1 的对角矩阵构成的子群，而所有与 $\Theta$ 的元素相联系的权空间都是一维的：
\[
\text{对一切 } \beta=\alpha_{i,i+1}\in\Theta,\qquad
\mathfrak{u}_\beta=\mathfrak{g}_\beta=\langle E_{i,i+1}\rangle.
\]
在 Lusztig 的全正性理论中，分裂实群的正元素半群由形如 $u_i(t)=I_n+tE_{i,i+1}=\exp(tE_{i,i+1})$（$t>0$）的元素生成。此时 $\mathfrak{u}_\beta=\mathfrak{g}_{\alpha_{i,i+1}}$ 中的尖锐凸锥就是 $E_{i,i+1}$ 的正数倍之锥。

% [Tutor Note: 这里给出全正性 $\leftrightarrow$ Θ-正性的对应关系：$\Theta=\Delta$ 时 $L_\Theta$ 是对角群，$\mathfrak{u}_\beta$ 的每个权空间一维，取“正数倍”即为尖锥；一般 $\Theta$ 时 $\mathfrak{u}_\beta$ 可能是高维不可约表示（Thm. 12.12），尖锥的存在性就不显然了——这正是 Guichard–Wienhard 分类定理（Thm. 12.16）的内容。]

\section{$\Sp_{2n}(\R)$ 中的 Θ-正性}

对管型 Hermitian 李群，存在另一种 Θ-正结构。我们针对实辛群 $\Sp_{2n}(\R)$ 引入这一结构，同时记住：对其他管型 Hermitian 李群，构造是完全类似的。

在这种情形下，Θ-正性的概念归结为 [69, 70] 中引入并研究的相应概念。取单根 $\Theta=\{\alpha_n\}$，其中 $\alpha_n$ 是 Dynkin 图末端、带有两个指向外侧箭头的那个根。此时只有一个权空间 $\mathfrak{u}_\beta$，即
\[
\mathfrak{u}_{\alpha_n}=\left\{\begin{pmatrix} 0 & B\\ 0 & 0\end{pmatrix}\ \Big|\ B\in\operatorname{Sym}_n(\R)\right\}.
\]
此外，$\mathfrak{l}_{\alpha_n}$ 与 $\mathfrak{p}_{\alpha_n}$ 分别与 11.1.2 节中 Levi 子群 $L$ 与抛物子群 $P$ 的李代数一致，并且可显式写为
\begin{align*}
\mathfrak{l}_{\alpha_n}&=\left\{\begin{pmatrix} A & 0\\ 0 & -A^t\end{pmatrix}\ \Big|\ A\in\operatorname{Mat}_n(\R)\right\},\\
\mathfrak{p}_{\alpha_n}&=\left\{\begin{pmatrix} A & B\\ 0 & -A^t\end{pmatrix}\ \Big|\ A\in\operatorname{Mat}_n(\R),\ B\in\operatorname{Sym}_n(\R)\right\}.
\end{align*}
存在一个开的尖锐凸锥

% [原文 p.138]
\[
\mathfrak{u}_{\alpha_n}^+=\left\{\begin{pmatrix} 0 & B\\ 0 & 0\end{pmatrix}\ \Big|\ B\in\operatorname{Sym}_n^+(\R)\right\}\subset\mathfrak{u}_{\alpha_n},
\]
它在 $L$ 的共轭作用下不变。于是它给出了 $\Sp_{2n}(\R)$ 上的一个正结构。注意，若 $n\ge 2$，该锥不是一维的，因此这个正结构与由 $\Sp_{2n}(\R)$ 是分裂实群这一事实所产生的全正结构不同。若 $n=1$，则 $\Sp_2(\R)=\SL_2(\R)$，此时 Θ-正性的概念与上一节讨论的全正性一致。

% [Tutor Note: “Dynkin 图末端带有两个指向外侧箭头”指 $C_n$ 型 Dynkin 图：$\alpha_1$ 为短根、$\alpha_n$ 为长根，双箭头由 $\alpha_{n-1}$ 指向 $\alpha_n$（在本书的约定下箭头指向长根）。$\alpha_n$ 对应的 $\mathfrak{u}_{\alpha_n}$ 是“最右上”的 $B$ 块；$\operatorname{Sym}_n^+(\R)$ 表示正定对称矩阵之锥，它在 $\GL_n(\R)$ 的合同作用下不变，而 $L$ 通过 $A\mapsto A B A^t$ 作用在 $B$ 上，故该锥不变。$n=1$ 时 $\operatorname{Sym}_1^+(\R)=\R_{>0}$ 是一维锥，恰回到全正性。]
